\subsection{在线性表示中的应用}

命题 13. 设 G 是连通李群，\(\rho\) 是 G 在完备可赋范空间 E 上的解析线性表示。令 \(\mathrm{E}_1, \mathrm{E}_2\) 为 E 的两个闭向量子空间，满足 \(\mathrm{E}_2 \subset \mathrm{E}_1\)。以下条件等价：

\begin{enumerate}
\def\labelenumi{(\roman{enumi})}
\item
  \(\rho(g)x \equiv x \pmod{\mathbf{E}_2}\)，对所有 \(g \in \mathbf{G}\) 和所有 \(x \in \mathbf{E}_1\)；
\item
  \(\mathrm{L}(\rho)(\mathrm{L}(\mathrm{G}))\) 把 \(E_{1}\) 映入 \(E_{2}\)。
\end{enumerate}

\[
\begin{array}{l l l} \rho (g) x \equiv x & (\mathrm{mod} \mathrm{E} _ {2}) & \text {for all \(g\in G\) and all \(x\in E_{1}\)} \\ \Leftrightarrow \rho (\exp a) x \equiv x & (\mathrm{mod} \mathrm{E} _ {2}) & \text {for all \(a\in L(G)\) and all \(x\in E_{1}\)} \\ \Leftrightarrow (\exp \mathrm{L} (\rho) a) x \equiv x & (\mathrm{mod} \mathrm{E} _ {2}) & \text {for all \(a\in L(G)\) and all \(x\in E_{1}\)}. \end{array}
\]

另一方面，若 \(u \in \mathcal{L}(\mathrm{E})\)，则

\[
\begin{array}{r l} & \exp (\lambda u) x \equiv x \pmod {\mathrm{E} _ {2}} \quad \text {(for all} \lambda \in \mathrm{K} \text {and all} x \in \mathrm{E} _ {1} \\ & \Leftrightarrow u (\mathrm{E} _ {1}) \subset \mathrm{E} _ {2} \end{array}
\]

从而命题得证。

推论 1. \(\mathbf{E}_1\) 在 \(\rho\) 下稳定的充要条件是 \(\mathbf{E}_1\) 在 \(\mathrm{L}(\rho)\) 下稳定。

在命题 13 中取 \(E_{1} = E_{2}\) 即可。

推论 2. 假设 \(\rho\) 是有限维的。\(\rho\) 为单的（分别为半单的）的充要条件是 \(\mathbf{L}(\rho)\) 为单的（分别为半单的）。

这由推论 1 得出。

推论 3. 设 \(x \in E\)。\(x\) 在 \(\rho(G)\) 下不变的充要条件是 \(x\) 被 \(L(\rho)(L(G))\) 消去（即，按 Chapter I, § 3, Definition 3 的意义，\(x\) 在 \(L(\rho)\) 下不变）。

在命题 13 中取 \(\mathbf{E}_1 = \mathbf{K}x\) 和 \(\mathbf{E}_2 = 0\) 即可。

推论 4. 设 \(\rho'\) 是 \(G\) 在完备可赋范空间 \(E'\) 上的另一个解析线性表示。令 \(T \in \mathcal{L}(E, E')\)。以下条件等价：

\begin{enumerate}
\def\labelenumi{(\roman{enumi})}
\item
  \(\mathrm{T}\rho (g) = \rho '(g)\mathrm{T}\)，对所有 \(g\in \mathbf{G}\)
\item
  \(\mathrm{TL}(\rho)(a) = \mathrm{L}(\rho')(a)\mathrm{T}\)，对所有 \(a\in \mathrm{L}(\mathrm{G})\)
\end{enumerate}

令 \(\sigma\) 为 \(G\) 在 \(\mathcal{L}(E, E')\) 上的线性表示，它由 \(\rho\) 和 \(\rho'\) 导出（\(\S 3\)，no. 11, Corollary 1 to Proposition 41）。条件 (i) 意味着 \(T\) 在 \(\sigma(G)\) 下不变。条件 (ii) 意味着 \(T\) 被 \(L(\sigma)(L(G))\) 消去。因此只需应用推论 3。

推论 5. 假设 \(\rho\) 和 \(\rho'\) 是有限维的。\(\rho\) 和 \(\rho'\) 等价的充要条件是 \(\mathbf{L}(\rho)\) 和 \(\mathbf{L}(\rho')\) 等价。

这是推论 4 的特殊情形。

推论 6. 假设 \(\mathbf{G}\) 是有限维的。令 \(t \in \mathrm{U}(\mathrm{G})\)。\(\mathbf{L}_t\) 为右不变（分别地，\(\mathbf{R}_t\) 为左不变）的充要条件是 \(t\) 属于 \(\mathrm{U}(\mathrm{G})\) 的中心。

\(\mathrm{L}_t\) 为右不变（分别地，\(\mathbf{R}_t\) 为左不变）的充要条件是 \(\varepsilon_g * t = t * \varepsilon_g\)，对所有 \(g \in G\)，也就是 \((\operatorname{Int} g)_* t = t\)。存在一个整数 \(n\)，使得 \(t \in \mathrm{U}_n(G)\)。由 § 3, no. 12 的 Corollary 3 和 Proposition 45，\((\operatorname{Int} g)_* t = t\) 对所有 \(g \in G\) 成立，当且仅当 \([a, t] = 0\) 对所有 \(a \in \mathrm{L}(G)\) 成立，也就是当且仅当 \(t\) 与 \(\mathrm{U}(G)\) 交换。

